\chapter{Introduction}
\label{ch:introduction}

This is a starter chapter to help you get going with the template.
Replace this content with your own writing. Below are examples of the
custom commands available in this template.


\section{Images}
\label{sec:images}

Use \verb|\insertimage| to include a figure with a caption and label:

\insertimage{placeholder.png}{A sample placeholder image}{fig:sample-image}{0.4\textwidth}

As shown in Figure~\ref{fig:sample-image}, images are automatically centered
and numbered.


\section{Tables}
\label{sec:tables}

Use \verb|\inserttable| to create tables:

\inserttable{A sample table}{tab:sample}{lcc}{
  \toprule
  \textbf{Item} & \textbf{Value} & \textbf{Status} \\
  \midrule
  Alpha & 42   & OK \\
  Beta  & 17   & OK \\
  Gamma & 99   & N/A \\
  \bottomrule
}

Refer to Table~\ref{tab:sample} using its label.


\section{Code Snippets}
\label{sec:code}

Insert code from an external file:

\insertcode{code/sample.py}{python}{A sample Python script}{lst:sample}

Or show specific lines:

\insertcodelines{code/sample.py}{python}{6}{9}{function get\_avatar}{lst:sample-lines}

\insertlinescodeclean{code/sample.py}{python}{11}{13}{API handler}{lst:handler}

\section{Callout Boxes}
\label{sec:callouts}

\note{This is an informational note. Use it to highlight important details.}

\warningbox{This is a warning box. Use it for caveats and potential pitfalls.}

\tipbox{This is a tip box. Use it for best practices and helpful hints.}


\section{Abbreviations}
\label{sec:abbreviations}

First mention of \abbr{api} will include a footnote with the full form.
Subsequent uses of \abbr{api} will just show the short form.

You can also use:
\begin{itemize}[leftmargin=2em]
  \item \verb|\abbrfull{api}|: \abbrfull{api}
  \item \verb|\abbrshort{api}|: \abbrshort{api}
  \item \verb|\abbrlong{api}|: \abbrlong{api}
\end{itemize}


\section{References}
\label{sec:references}

Cite sources with \verb|\insertref{key}|, e.g.\
\insertref{knuth1984} and \insertref{lamport1994}.
